\documentclass[11pt]{article}
\usepackage[letterpaper,margin=0.68in,headheight=15pt]{geometry}
\usepackage[T1]{fontenc}
\usepackage{lmodern}
\usepackage{amsmath,amssymb}
\usepackage{array,booktabs,enumitem,multicol}
\usepackage{xcolor}
\usepackage{tikz}
\usepackage{fancyhdr}
\usepackage{microtype}
\setlength{\parindent}{0pt}
\setlength{\parskip}{5pt}
\definecolor{olive}{HTML}{555B35}
\definecolor{gold}{HTML}{B18A22}
\definecolor{soft}{HTML}{F4F1DE}
\definecolor{bluegray}{HTML}{E8EEF2}
\pagestyle{fancy}
\fancyhf{}
\fancyhead[L]{\textcolor{olive}{\small MATH\& 141: Precalculus I}}
\fancyhead[R]{\textcolor{gold}{\small Fall 2026}}
\fancyfoot[C]{\thepage}
\renewcommand{\headrulewidth}{0.45pt}
\renewcommand{\footrulewidth}{0pt}
\setlist[enumerate]{leftmargin=*,itemsep=0.08in,topsep=0.05in}
\newcommand{\assessmenttitle}[2]{%
  {\Large\bfseries Module #1: #2}\par\vspace{3pt}
}
\newcommand{\studentline}{%
\noindent\textbf{Name:}\hspace{2.6in}\textbf{Date:}\hspace{1.15in}\par\vspace{7pt}}
\newcommand{\directions}[1]{\textit{#1}\par\vspace{4pt}}
\newcommand{\answerblank}[1]{\vspace{#1}}
\newcommand{\blankgrid}[4]{%
\begin{center}
\begin{tikzpicture}[x=0.75cm,y=0.75cm]
  \draw[step=1,gray!28,very thin] (#1,#3) grid (#2,#4);
  \draw[->,thick] (#1,0)--(#2,0) node[right] {$x$};
  \draw[->,thick] (0,#3)--(0,#4) node[above] {$y$};
  \foreach \x in {#1,...,#2}{\ifnum\x=0\else\draw (\x,0.10)--(\x,-0.10) node[below=2pt,font=\scriptsize]{\x};\fi}
  \foreach \y in {#3,...,#4}{\ifnum\y=0\else\draw (0.10,\y)--(-0.10,\y) node[left=2pt,font=\scriptsize]{\y};\fi}
\end{tikzpicture}
\end{center}}
\newcommand{\smallgrid}[4]{%
\begin{tikzpicture}[x=0.50cm,y=0.50cm]
  \draw[step=1,gray!28,very thin] (#1,#3) grid (#2,#4);
  \draw[->,thick] (#1,0)--(#2,0) node[right] {$x$};
  \draw[->,thick] (0,#3)--(0,#4) node[above] {$y$};
\end{tikzpicture}}
\begin{document}

% MODULE 6 PREQUIZ
\assessmenttitle{6}{Exponential Functions -- Prequiz}
\studentline
\directions{Four short questions. Interpret the initial value and repeated multiplicative change.}
\begin{enumerate}
\item For $A(t)=400(1.06)^t$, state the initial value and the percent growth rate.
\answerblank{0.48in}
\item A population is modeled by $P(t)=250\,2^{t/5}$. Find $P(10)$.
\answerblank{0.52in}
\item A substance begins with $960$ grams and has a half-life of $4$ hours. How much remains after $12$ hours?
\answerblank{0.58in}
\item A continuously growing quantity is modeled by $B(t)=500e^{0.03t}$. Find $B(5)$ to three decimal places.
\answerblank{0.62in}
\end{enumerate}
\newpage

% MODULE 6 CHECKIN PAGE 1
\assessmenttitle{6}{Exponential Functions -- Weekly Check-In}
\studentline
\directions{Three questions. Show enough work to make the repeated multiplicative structure visible.}
\textbf{1. Reading an exponential model.}

A population is modeled by
\[
P(t)=450(1.08)^t,
\]
where $t$ is measured in years.
\begin{enumerate}[label=(\alph*),leftmargin=0.35in]
\item State the initial population and the percent growth rate.
\answerblank{0.42in}
\item Find $P(6)$ to three decimal places.
\answerblank{0.62in}
\end{enumerate}

\textbf{2. Half-life.}

A medication begins with $640$ mg in the body and has a half-life of $6$ hours.
\begin{enumerate}[label=(\alph*),leftmargin=0.35in]
\item Write an exponential model for the amount $A(t)$ after $t$ hours.
\answerblank{0.48in}
\item Find $A(9)$ to three decimal places and $A(18)$ exactly.
\answerblank{0.75in}
\end{enumerate}
\newpage

% MODULE 6 CHECKIN PAGE 2
\assessmenttitle{6}{Exponential Functions -- Weekly Check-In, continued}
\studentline
\textbf{3. Continuous exponential growth.}

A culture begins with $300$ cells and is modeled by
\[
C(t)=300e^{0.08t},
\]
where $t$ is measured in hours.
\begin{enumerate}[label=(\alph*),leftmargin=0.35in]
\item Explain what the numbers $300$ and $0.08$ tell you about the model.
\answerblank{0.85in}
\item Find $C(5)$ to three decimal places.
\answerblank{0.85in}
\item Is the quantity growing or decaying? Explain how the formula tells you.
\answerblank{0.85in}
\end{enumerate}

\end{document}
